\section*{Bab \ref{ch:g11:work-of-force} --- Usaha sebuah Gaya}

\begin{solution}{exo:g11:work-of-force:1}
$W = 40 \times 50 \times \cos\theta$: $+\qty{2000}{J}$;
$+\qty{1000}{J}$; $0$; $-\qty{1000}{J}$.
\end{solution}

\begin{solution}{exo:g11:work-of-force:2}
$mgh = 12 \times 9.81 \times 2.5 = \qty{294}{J}$: ketika diturunkan,
$+\qty{294}{J}$ (\omterm{def:g10:universal-gravitation:weight}{beratnya} searah turunnya); ketika dinaikkan,
$-\qty{294}{J}$.
\end{solution}

\begin{solution}{exo:g11:work-of-force:3}
(a) \omterm{def:g11:work-of-force:sign}{pendorong}; (b) \omterm{def:g11:work-of-force:sign}{penghambat}; (c) nol (tegak lurus); (d) \omterm{def:g11:work-of-force:sign}{penghambat};
(e) nol --- tegangan talinya selalu tegak lurus terhadap gerak
bandulnya.
\end{solution}

\begin{solution}{exo:g11:work-of-force:4}
Laju tetap: \omterm{def:g10:inertia:force}{gaya} pengangkatnya $= mg = \qty{5.89e3}{N}$, sehingga
$W_{\text{angkat}} = 5886 \times 25 \approx +\qty{147}{kJ}$ dan
$W(\vect P) = -\qty{147}{kJ}$. Jumlahnya nol: pada kecepatan tetap
\omterm{def:g10:inertia:force}{gayanya} setimbang (\cref{ch:g11:forces-and-motion}), sehingga \omterm{def:g11:work-of-force:work}{usahanya}
saling menghapus.
\end{solution}

\begin{solution}{exo:g11:work-of-force:5}
$P = \num{9.0e4}/60 = \qty{1.5}{kW}$; $1500/736 \approx \qty{2.0}{hp}$.
\end{solution}

\begin{solution}{exo:g11:work-of-force:6}
Talinya: $120 \times 200 \times \cos\qty{30}{\degree} \approx
+\qty{20.8}{kJ}$; \omterm{def:g10:inertia:inventory}{gesekannya}: $-45 \times 200 = -\qty{9.0}{kJ}$; \omterm{def:g10:universal-gravitation:weight}{berat}
dan \omterm{def:g10:inertia:inventory}{gaya normalnya}, yang tegak lurus: $0$. Seluruhnya
$+\qty{11.8}{kJ}$: kereta luncur itu bertambah cepat.
\end{solution}

\begin{solution}{exo:g11:work-of-force:7}
$v = 25/3.6 = \qty{6.94}{m/s}$; $P = Fv = 18 \times 6.94 \approx
\qty{125}{W}$. Dalam satu jam, $W = P\,\Delta t = 125 \times 3600 =
\qty{450}{kJ}$ (setara dengan $18 \times \qty{25000}{m}$).
\end{solution}

\begin{solution}{exo:g11:work-of-force:8}
\omterm{def:g10:universal-gravitation:weight}{Beratnya}, pada kedua rute: $+65 \times 9.81 \times 120 \approx
+\qty{76.5}{kJ}$ (hanya selisih ketinggiannya). \omterm{def:g10:inertia:inventory}{Gesekannya}:
$-30 \times 800 =
-\qty{24}{kJ}$ pada lerengnya,
$-30 \times 1500 = -\qty{45}{kJ}$ pada jalurnya. Hanya \omterm{def:g10:inertia:inventory}{gesekannya} yang
peduli pada rutenya.
\end{solution}

\begin{solution}{exo:g11:work-of-force:9}
Kenaikannya $h = 4.0\sin\qty{25}{\degree} = \qty{1.69}{m}$:
$W(\vect P) = -20 \times 9.81 \times 1.69 \approx -\qty{332}{J}$;
$W_f = -30 \times 4.0 = -\qty{120}{J}$. Karena lajunya tetap, \omterm{def:g11:work-of-force:work}{usaha}
seluruhnya nol, sehingga $W_{\text{dorong}} = +\qty{452}{J}$ dan
$F = 452/4.0 \approx \qty{113}{N}$.
\end{solution}

\begin{solution}{exo:g11:work-of-force:10}
$v = 130/3.6 = \qty{36.1}{m/s}$; karena lajunya tetap,
$F = P/v = \num{5.0e4}/36.1 \approx \qty{1.4}{kN}$.
\end{solution}

\begin{solution}{exo:g11:work-of-force:11}
$P = Fv\cos\theta = 800 \times 1.0 \times \cos\qty{20}{\degree}
\approx \qty{752}{W}$; dalam \omterm{def:g11:work-of-force:power}{daya kuda}, $752/736 \approx \num{1.02}$
--- hampir persis yang sanggup ditahan seekor kuda yang kuat, yaitu
perbandingan yang dipakai Watt ketika membangun \omterm{def:g10:orders-of-magnitude:unit}{satuannya}.
\end{solution}

\begin{solution}{exo:g11:work-of-force:12}
Panjangnya $L = h/\text{tanjakan}$: \qty{1500}{m} dan \qty{5000}{m}.
\omterm{def:g10:inertia:force}{Gayanya} $F = mg \times \text{tanjakan}$:
$1200 \times 9.81 \times 0.20 \approx
\qty{2.35}{kN}$ dan
$\qty{0.71}{kN}$. \omterm{def:g11:work-of-force:work}{Usahanya} $FL$: $2354 \times 1500
\approx \qty{3.53}{MJ}$ dan $706 \times 5000 \approx \qty{3.53}{MJ}$ ---
keduanya sama dengan $mgh$: jalannya menukar \omterm{def:g10:inertia:force}{gaya} dengan panjang, tidak
pernah dengan \omterm{def:g11:work-of-force:work}{usaha}.
\end{solution}

\begin{solution}{exo:g11:work-of-force:13}
$F = mg = \qty{2.45e3}{N}$;
$W = 2453 \times 15 \approx \qty{36.8}{kJ}$;
$\Delta t = 15/0.40 = \qty{37.5}{s}$;
$P = 36800/37.5 \approx \qty{981}{W} = 2453 \times 0.40$ ($Fv$ sesuai);
\omterm{prop:g11:circuits-and-power:power}{daya listriknya}: $981/0.70 \approx \qty{1.4}{kW}$.
\end{solution}

\begin{solution}{exo:g11:work-of-force:14}
Diagonalnya: $-140 \times 6.0 = -\qty{840}{J}$; dindingnya:
$-140 \times 8.4 \approx -\qty{1.18}{kJ}$; pulang perginya:
$-140 \times 12 = -\qty{1.68}{kJ}$. \omterm{def:g10:universal-gravitation:weight}{Beratnya} melakukan nol pada
perjalanan pulang pergi itu (tidak ada selisih ketinggian). \omterm{def:g11:work-of-force:work}{Usaha} oleh
\omterm{def:g10:inertia:inventory}{gesekan} bergantung pada \omterm{def:g10:relative-motion:trajectory}{lintasannya} dan tidak pernah mengembalikan;
\omterm{def:g11:work-of-force:work}{usaha} oleh \omterm{def:g10:universal-gravitation:weight}{berat} hanya bergantung pada kedua ujungnya
(\cref{rem:g11:work-of-force:divide}).
\end{solution}

\begin{solution}{exo:g11:work-of-force:15}
Melawan gravitasi: $80 \times 9.81 \times 600 \approx \qty{471}{kJ}$;
melawan \omterm{def:g10:inertia:inventory}{gesekan}: $14 \times 12000 = \qty{168}{kJ}$; seluruhnya
$\qty{639}{kJ}$. $v = \qty{4.17}{m/s}$, sehingga
$\Delta t = 12000/4.17 = \qty{2880}{s} =
\qty{48}{min}$ dan
$P = \num{6.39e5}/2880 \approx \qty{222}{W} \approx
\qty{0.30}{hp}$ ---
seorang amatir yang kuat, sepertiga ekor kuda.
\end{solution}

\begin{solution}{pb:g11:work-of-force:1}
\textbf{1.} $W = 110 \times 8.0 = +\qty{880}{J}$.

\textbf{2.} \omterm{def:g10:inertia:inventory}{Gesekannya}: $-95 \times 8.0 = -\qty{760}{J}$. \omterm{def:g10:universal-gravitation:weight}{Berat} dan
\omterm{def:g10:inertia:inventory}{gaya normalnya} tegak lurus geraknya: masing-masing nol.

\textbf{3.} Seluruhnya $+\qty{120}{J}$: \omterm{def:g11:work-of-force:sign}{usaha pendorong} bersihnya,
sehingga kotak itu bertambah cepat. Dorongan \qty{95}{N} mengimbangi
\omterm{def:g10:inertia:inventory}{gesekannya} --- \omterm{def:g11:forces-and-motion:netforce}{gaya totalnya} nol, lajunya tetap, dan \omterm{def:g11:work-of-force:work}{usaha} seluruhnya
nol.

\textbf{4.} Diagonalnya: $-95 \times 10.0 = -\qty{950}{J}$;
dindingnya: $-95 \times 14.0 = -\qty{1.33}{kJ}$: \omterm{def:g10:inertia:inventory}{gesekan} menagih tiap
\omterm{def:g10:orders-of-magnitude:unit}{meter}.

\textbf{5.} \omterm{def:g10:inertia:inventory}{Gesekannya}: $-95 \times 20.0 = -\qty{1.9}{kJ}$; \omterm{def:g10:universal-gravitation:weight}{beratnya}:
$0$ (tidak ada selisih ketinggian). Hanya \omterm{def:g10:universal-gravitation:weight}{beratnya} yang mengembalikan.

\textbf{6.} $W(\vect P) = -30 \times 9.81 \times 9.0 =
-\qty{2.65}{kJ}$; regunya sekurang-kurangnya harus memasok
$+\qty{2.65}{kJ}$.

\textbf{7.} $(70 + 30) \times 9.81 \times 9.0 = \qty{8.83}{kJ}$;
bagian milik kotaknya $2.65/8.83 \approx 30\%$ --- memanggul sebagian
besar berarti mengangkut diri sendiri.

\textbf{8.} $\sin\alpha = 9.0/30 = 0.30$:
$F = mg\sin\alpha + f = 88.3 + 40 = \qty{128}{N}$;
$W = 128 \times 30 = \qty{3.85}{kJ}$ ($= 2.65 + 1.2$).

\textbf{9.} Pengerekan tegak lurus: $F = mg = \qty{294}{N}$,
$W = \qty{2.65}{kJ}$. Bidang miringnya membelikan \omterm{def:g10:inertia:force}{gaya} yang lebih kecil
($128$ berbanding \qty{294}{N}) dengan harga \omterm{def:g11:work-of-force:work}{usaha} yang lebih besar:
retribusi \omterm{def:g10:inertia:inventory}{gesekan} sebesar \qty{1.2}{kJ}.

\textbf{10.} $\Delta t = 9.0/0.50 = \qty{18}{s}$;
$P = 2649/18 \approx \qty{147}{W}$.

\textbf{11.} $P = Fv = 294 \times 0.50 = \qty{147}{W}$ --- sama.
\omterm{prop:g11:circuits-and-power:power}{Daya listriknya}: $147/0.60 \approx \qty{245}{W}$.

\textbf{12.} $1000 \times 0.24 = \qty{240}{m}$ dan
$8000 \times 0.030 = \qty{240}{m}$: puncak yang sama.

\textbf{13.} $F = mg \times \text{tanjakan}$: pada jalurnya
$3500 \times 9.81 \times 0.24 \approx \qty{8.2}{kN}$; pada jalannya
$\approx \qty{1.0}{kN}$.

\textbf{14.} $8240 \times 1000 \approx \qty{8.24}{MJ}$ dan
$1030 \times 8000 \approx \qty{8.24}{MJ}$ --- keduanya
$= mgh = 3500
\times 9.81 \times 240$.

\textbf{15.} \omterm{def:g11:circuits-and-power:resistance}{Hambatan} gelindingnya:
$400 \times 1000 = \qty{0.40}{MJ}$ berbanding
$400 \times 8000 = \qty{3.2}{MJ}$: kenyamanan jalan yang landai itu
berharga \qty{2.8}{MJ} \omterm{def:g10:inertia:inventory}{gesekan} tambahan.

\textbf{16.} $v = \qty{10}{m/s}$. Jalannya:
$(1030 + 400) \times 10 = \qty{14.3}{kW} \approx \qty{19}{hp}$;
jalurnya: $(8240 + 400) \times 10 = \qty{86.4}{kW} \approx \qty{117}{hp}$ --- di luar \omterm{def:g11:fundamental-interactions:strong}{jangkauan} sebuah van bermuatan. Jalan
berkelok agar mesin (dan ban) yang sederhana sanggup mendaki: \omterm{def:g11:work-of-force:work}{usaha}
yang sama, disebar \omterm{def:g11:lenses-and-eye:lens}{tipis}.

\textbf{17.} Tiap kotaknya: $880 + 2649 = \qty{3.53}{kJ}$; untuk $60$
kotak, $60 \times 3529 \approx \qty{212}{kJ}$.

\textbf{18.} Tangganya: $20 \times 70 \times 9.81 \times 9.0 \approx
\qty{124}{kJ}$. Jumlah seluruhnya untuk hari itu
$\approx \qty{336}{kJ} = \qty{0.34}{MJ}$.

\textbf{19.} $P = \num{3.36e5}/28800 \approx \qty{12}{W}$. Ketelnya:
$\num{3.36e5}/2000 = \qty{168}{s}$ --- kurang dari tiga menit untuk
kerja keras sehari penuh.

\textbf{20.} Joule membayar \omterm{def:g10:inertia:force}{gaya} yang diserahkan \emph{sepanjang}
sebuah gerak --- dan tidak membayar apa pun untuk dinding yang didorong
atau kotak yang ditahan, seletih apa pun rasanya (otot membakar \omterm{def:g10:energy-conservation:forms}{energi
kimia} bahkan ketika \omterm{def:g11:work-of-force:work}{usaha} mekanisnya nol). Karena keluaran \omterm{def:g11:work-of-force:power}{daya} manusia
hanya beberapa puluh \omterm{def:g11:work-of-force:power}{watt}, para pengangkut barang membiarkan mesin
membentuk ulang \omterm{def:g10:inertia:force}{gayanya} --- derek gulung, bidang miring, jalan berkelok
--- sementara tagihan $mgh$ tidak pernah berubah.
\end{solution}
